%! Author = chtompki
%! Date = 1/14/26
% Example LaTeX document for GP111 - note % sign indicates a comment
\documentclass[12pt]{amsart}
% Default margins are too wide all the way around. I reset them here
\setlength{\hoffset}{-.75in}
\setlength{\textwidth}{6.5in}
\setlength{\voffset}{-.5in}
\setlength{\textheight}{9.0in}
\usepackage{amsmath}
\usepackage{amssymb}
\usepackage{amsthm}
\usepackage{pgfplots}
\usepackage{enumerate}
\usepackage{tikz}
\usetikzlibrary{shadows,arrows,arrows.meta,shapes,positioning,calc}% Required for the Circle and Triangle arrow tips
\usepackage{setspace}
\usepackage{siunitx}

\theoremstyle{definition}
\newtheorem{definition}{Definition}

\theoremstyle{question}
\newtheorem{question}{Question}

\title{Seminar Summary:\\February 24, 2026 - Laura Fix\\Patient-specific modeling of cardiovascular and respiratory dynamics}
\author{Rob Tompkins}
\date{\today}


\begin{document}
    \maketitle
    \date{\today}
    \begin{onehalfspace}

        \section{Cardiorespiratory compartmental model}
        The utilization of a mathematical model in regard to modeling the dynamics of blood-flow through the body
        is informative and understood to be fairly predictive of the actual dynamics of the physiological system
        at hand.
        A standard in the field of fluid dynamics for modeling flow through a network of pipes, pumps, check valves,
        capacitor tanks.
        Ohms law in fluid dynamics follows as
        \begin{displaymath}
            P = Q\cdot R
        \end{displaymath}
        where $P$ is pressure head at a spot along the pipe, $Q$ is the flow rate, and $R$ is the resistence of the
        pipe.
        Now this is a description of the model that we were presented and is an ODE model.
        Generally the system can be modeled with ODEs or PDEs.

        We were given two projects that were being actively worked on: prediction of dynamic cerebrovascular resistance and
        patient-specific modeling of respiratory mechanics in preterm infants (which seems to have new material to work on).

        We then briefly went over the mechanics of orthostatics and how sudden shifts in body position accommodate for
        blood quickly moving throughout the system with reference to gravity such that cranial blood pressure can drop
        below functional levels.
        We then delved into the general physiology of the circulatory system as well as the autonomic nervous system,
        with the sympathetic and parasympathetic controls that it offers.

        The model was set up as eight distinct compartments over the arterial and veinous vascular systems,
        each of which had four complementary compartments: the head, the heart, the thorax, and systemic tissue (the
        remainder catchall). In a rough sketch the model was given as
        \begin{displaymath}
                \frac{dV}{dt} = q_{in} - q_{out},\quad\text{where}\quad q = \frac{p_{in} - p_{out}}{R}
        \end{displaymath}
        and
        \begin{displaymath}
            V - V_{unstr} = Cp\quad \Rightarrow\quad \frac{dp}{dt} = \frac{1}{C}\left(q_{in}-q_{out}-p\frac{dC}{dt}\right)
        \end{displaymath}

        The model was then run against normal respiration followed the administration of a breathing mask
        giving 5\% $CO_2$, and plots were shown.
        It was shown that cranial resistence hops up after re-breathing the 5\% $CO_2$ for some time.
        There seem to be future research opportunities here regarding applying the model to populations with
        compromised cerebrovascular reactivity.
        Another thought was to potentially include modern sensitivity analysis and parameter identifiability
        techniques.

        \section{Patient-specific modeling of respiratory mechanics in the preterm infant}
        When considering the 1st-percentile in infant birth weight we are faced with a mortality rate
        of 60\%.
        The issue seems to stem from high chest wall compliance due to an undermineralized rib cage, and cascades
        of comorbidities can quickly arise necessitating ventilation.

        Now a paper by Abbasi and Bhutani (1990) was presented giving data with clinical pressures, flows, and volumes.
        We compared their data with neurally adjusted ventilatory assist data and they aligned well.
        (\emph{Note.} neurally adjusted ventilatory assist is a mechanism for ventilation that relies upon a sensor catheter
        placed in the esophagus of the patient to get electrical conductivity from the autonomic nervous system).

        Similarly to before we fall back on an electromagnetic fluid dynamics analogue for modeling the alveoli of the
        neonate needing care, however this time the mechanics of the fluid are complex enough to yield a dynamical system
        as opposed to a solvable set of ODEs.

        Using Matlab for modeling the system the plots shown indicated that the model did indeed align with dynamical system
        we created.
        Then a fairly wide array of simulated experiments were run through the model like the patient on a CPAP machine,
        what happens with an injured lung, and what happens under anesthesia.

        We went over a little more about sensitivity analysis, but then we wrapped things up.


    \end{onehalfspace}
\end{document}
